% !TEX TS-program = lualatex
%\documentclass[handout]{beamer}
\documentclass{beamer}
\input{../common/misomva}
\newcommand{\misoclasstinytitleslide}{Partially based on material by:  Ulrike von Luxburg, \\[-1em] Gary Miller, Doyle \& Schnell, Daniel Spielman}
\newcommand{\misoclassdate}{}
\usetikzlibrary{graphs}

\begin{document}
% Topic-specific title slide

\newcommand{\topictitle}{Normalized Laplacians}

\newcommand{\topicsubtitle}{$L_{\mathrm{sym}}$ and $L_{\mathrm{rw}}$}

\MakeTopicSlide

\begin{frame}
\frametitle{Smoothness of the Function and Laplacian}
\begin{block}{}
$$
S_G(\bff)  = \bff\transpose\bL\bff =
 \bff\transpose\bQ\bLambda\bQ\transpose\bff = 
\balpha\transpose\bLambda\balpha = 
\|\balpha\|^2_{\bLambda}=\sum_{i=1}^\nodes\lambda_i\alpha_i^2
$$
 \end{block}

Eigenvectors are graph functions too!\\
\smallskip \pause
\questionbox{What is the smoothness of an eigenvector?}\\
\smallskip \pause
Spectral coordinates of eigenvector $\bv_k$: $\bQ\transpose\bv_k =  \pause \be_k$
\pause
\begin{block}{}
$$
S_G(\bv_k)\! = \! \bv_k\transpose\bL\bv_k \!=\!
 \bv_k\transpose\bQ\bLambda\bQ\transpose\bv_k = 
\be_k\transpose\bLambda\be_k = 
\|\be_k\|^2_{\bLambda}=\sum_{i=1}^\nodes\lambda_i(\be_k)_i^2 \pause = \lambda_k
$$
 \end{block}
\pause
The smoothness of $k$-th eigenvector is the $k$-th eigenvalue.


\end{frame}
%%%%%%%%%%%%%%%%%%%%%%%%%%%%%%%%%%%%%%%%%%%%%%%%%%%%%%%%%%%%%

\begin{frame}
\frametitle{Laplacian of the Complete Graph $K_\nodes$}

What is the eigenspectrum of $\bL_{K_\nodes}$?

\begin{columns}[c]
\column{0.8in}
\tikz \graph[nodes={circle,draw},edges={circle connection bar, fill}] 
{ subgraph K_n [n=5, counterclockwise] };

\column{2in}
\begin{center}
 {\scriptsize
\specialcell{
\begin{math}
\bL_{K_\nodes} = \left(
\begin{array}{ccccc}
\nodes -1 & -1 & -1 & -1 & -1 \\
 -1 & \nodes-1 & -1 & -1 & -1 \\
 -1 & -1 & \nodes-1 & -1 & -1 \\
 -1 & -1 & -1 & \nodes-1 & -1 \\
 -1 & -1 & -1 & -1 & \nodes-1 \\
\end{array}
\right)
\end{math}
}
}

\end{center}
\end{columns}

\pause
\bigskip

From before: we know that $(0, \bOne_\nodes)$ is an eigenpair.

\begin{block}{}

If $\bv \ne 0_\nodes$ and $\bv \perp \bOne_\nodes \implies \sum_i \bv_i = 0$. 
\pause
To get the other eigenvalues, we compute $(\bL_{K_\nodes}\bv)_1$ and divide by $\bv_1$ (wlog $\bv_1 \ne 0$). 
\pause
\vspace{-1em}
$$ 
(\bL_{K_\nodes}\bv)_1 = (\nodes-1)\bv_1 -  \sum_{i=2}^\nodes \bv_i = \nodes \bv_1. 
$$
\end{block}
\pause
\begin{flushright}
What are the remaining eigenvalues/vectors? 
\vspace{0.2em}
\pause
\onslide<handout:0>{\notecol{\tiny \\ Answer: $\nodes-1$ eigenvectors $\perp \bOne_\nodes$ for eigenvalue $\nodes$ with multiplicity $\nodes-1$.\\}}
\vspace{-0.5em}
%\gcol{\tiny  Question: What changes for \textbf{weighted} complete graphs?}
\end{flushright}
\end{frame}


%%%%%%%%%%%%%%%%%%%%%%%%%%%%%%%%%%%%%%%%%%%%%%%%%%%%%%%%%%%%%
\begin{frame}
\frametitle{Normalized Laplacians}

\begin{align*}
\onslide<1->{ \bL_{un} &= \bD - \bW }  \\ 
\onslide<2->{ \bL_{sym} &= \bD^{-1/2}\bL\bD^{-1/2} = \bI - \bD^{-1/2}\bW\bD^{-1/2} } \\ 
\onslide<3->{ \bL_{rw} &= \bD^{-1}\bL = \bI- \bD^{-1}\bW }  
\end{align*}


\onslide<4->{
\begin{block}{}
$$\bff\transpose\bL_{sym}\bff = \frac{1}{2}\sum_{i,j\leq \nodes}w_{i,j}%
\left(\frac{f_i}{\sqrt{d_i}}-\frac{f_j}{\sqrt{d_j}}\right)^2$$
\end{block}
}

\onslide<5->{
\begin{block}{}
$(\lambda,\bu)$ is an eigenpair for $\bL_{rw}$ iff 
$(\lambda,\bD^{1/2}\bu)$ is an eigenpair for $\bL_{sym}$
\end{block}\pause
}


\end{frame}
%%%%%%%%%%%%%%%%%%%%%%%%%%%%%%%%%%%%%%%%%%%%%%%%%%%%%%%%%%%%%



\MakeLastSlide
\end{document}
